\documentclass[12pt]{article}

\usepackage{sbc-template}
\usepackage{graphicx,url}
\usepackage[utf8]{inputenc}
\usepackage[brazil]{babel}
%\usepackage[latin1]{inputenc}  
\usepackage{listings}
\usepackage{xcolor}

% 1. Definindo as cores que você quer usar
\definecolor{codegreen}{rgb}{0,0.6,0}
\definecolor{codegray}{rgb}{0.5,0.5,0.5}
\definecolor{codepurple}{rgb}{0.58,0,0.82}
\definecolor{backcolour}{rgb}{0.95,0.95,0.92} % Cor de fundo suave

% 2. Configurando o estilo do código
\lstdefinestyle{mestilo}{
    backgroundcolor=\color{backcolour},   
    commentstyle=\color{codegreen}\itshape, % Comentários em verde e itálico
    keywordstyle=\color{codepurple}\bfseries,% Palavras-chave em roxo e negrito
    numberstyle=\tiny\color{codegray},      % Números de linha em cinza
    stringstyle=\color{codegreen},          % Strings em verde
    basicstyle=\ttfamily\footnotesize,      % Fonte padrão do código
    breakatwhitespace=false,         
    breaklines=true,                   % Quebra de linha automática
    captionpos=b,                      
    keepspaces=true,                   
    numbers=left,                      % Números de linha à esquerda
    numbersep=5pt,                     
    showspaces=false,                  
    showstringspaces=false,
    showtabs=false,                    
    tabsize=2,
    frame=single                       % Borda ao redor do código
}
     
\sloppy

\title{Classificação de Imagens de Plantas Doentes utilizando Visão Computacional}

\author{Henrique Mendonça Castelar Campos\inst{1}}


\address{Pontifícia Universidade Católica de Minas Gerais}

\begin{document} 

\maketitle

\begin{abstract}

In agriculture one of the challenges consists of the detection and isolation of diseased plants. Through computer vision, this process can be automated, reducing the losses of sick crops. In this work, have been developed Convolutional Neural Networks (CNNs) machine learning models to identify and detect sick plants. For this end, the following models have been developed: ResNet50, EfficientNetV2S and GoogleNet.
  
\end{abstract}
     
\begin{resumo} 

Na agricultura um dos desafios encontrados consiste na detecção e isolamento de plantas doentes. Através da visão computacional, esse processo pode ser automatizado, podendo reduzir as perdas resultantes de culturas doentes. Neste trabalho, foram desenvolvidos modelos de aprendizado de máquina utilizando Redes Neurais Convolucionais (CNNs) visando a detecção de plantas doentes. Para isso foram desenvolvidos os seguintes modelos: ResNet50, EfficientNetV2S e GoogleNet.
  
\end{resumo}

\section{Introdução}

\paragraph{}Um dos desafios encontrados na agricultura moderna consiste na identificação, detecção e isolamento de culturas de plantas doentes, tendo em vista que a doença pode se espalhar, resultando na contaminação de toda a produção agrícola. Dessa forma, a detecção automatizada de plantas doentes pode resultar no aumento na quantidade e na qualidade da produção agrícola, portanto reduzindo as perdas.

\paragraph{}Este trabalho tem como objetivo o desenvolvimento de um modelo de rede neural convolucional capaz de utilizar a visão computacional para identificar e classificar doenças em plantas. Dessa forma, agricultores poderão utilizar esse modelo para a automatização de identificação dessas doenças, impedindo que as suas culturas sofram perdas.

\section{Trabalhos Relacionados}

\paragraph{}Em \cite{10393095}, foi proposto o desenvolvimento de uma plataforma chamada “FarmEasy” na qual possui capacidade de detectar e classificar doenças de plantas. Para isso, foi utilizado o dataset New Plant Disease Dataset, que possuí 87867 imagens (70295 imagens de teste e 17572 imagens de treino) associadas a doenças de plantas. Esse dataset foi utilizado para treinar e testar os modelos VGG16, Inception v3 e ResNet50.

\paragraph{}Bahaa S. Hamed, Mahmoud M. Hussein, Afaf M. Mousa em \cite{article} propuseram o desenvolvimento de uma rede neural convolucional EfficientNetV2S para a classificação de doenças em plantas. Para isso, eles também utilizaram o dataset New Plant Disease Dataset no qual aplicaram uma série de modificações visando a tornar o modelo resiliente a ruídos.

\section{Metodologia}

\subsection{Dataset}

\paragraph{}Para o desenvolvimento deste trabalho, foi utilizado o Plant Village Dataset, o qual possuí 38 classes de doenças em plantas, incluindo classes de plantas saudáveis. A escolha deste dataset se deu pelo fato dele ser um dataset utilizado para a classificação de doenças de plantas.

\subsection{Pré-Processamento}

\paragraph{}Visando a facilitar o treinamento do modelo, foi gerado uma tabela no formato de arquivo CSV contendo os seguintes atributos: os nomes dos arquivos de imagem; a versão desse arquivo: color, grayscale e segmented; a classe; o tipo de planta; e o tipo de doença. Através dessa tabela, foi possível realizar a separação dos conjuntos de treino e teste que foram utilizados no treinamento dos modelos. A separação do conjunto de treino e de teste foi feita pela função ‘train\_test\_split’, à qual foi definido como critério de separação a margem de 20\% para o conjunto de teste, além do estado randômico com o valor 1024.

\subsection{Modelos}

\paragraph{}Neste trabalho foram utilizados os seguintes modelos: ResNet50, EfficientNetV2S e GoogleNet. A escolha dos modelos ResNet50 e EfficientNet se deu pelo fato de nos trabalhos relacionados analisados, esse modelos foram utilizados para o treinamento da rede neural convolucional. Dessa forma, o GoogleNet foi escolhido para ser o modelo a ser comparado com o EfficientNet e o ResNet.

\paragraph{}Um aspecto fundamental no desenvolvimento dos modelos foi o uso da biblioteca Tensorflow. Através dela foi possível utilizar os modelos pré-existentes ResNet50 e EfficientNetV2S. Dessa forma, não foi necessário implementar manualmente esses modelos.

\paragraph{}Em relação aos valores iniciais para os pesos dos modelos, foram utilizados os valores da ImageNet. Dessa forma, através do aprendizado por transferência, foi possível economizar tempo no treinamento dos modelos.

\lstinputlisting[
    breaklines=true,
    language={Python},
    caption={Código fonte da construção do modelo ResNet 50},
    label={codigo:resnet50},
    style=mestilo
]{Codigo/resnet50.py}

\lstinputlisting[
    breaklines=true,
    language={Python},
    caption={Código fonte da construção do modelo EfficientNetV2S},
    label={codigo:efficientnetv2s},
    style=mestilo
]{Codigo/efficientnetv2s.py}

\paragraph{}Em relação ao modelo GoogleNet, este teve de ser implementado manualmente, tendo em vista que o Tensotflow não possuí este modelo pronto. Dessa forma, foram desenvolvidas as seguintes funções: inception\_module, na qual realiza parte da convoluções realizadas pelo modelo;  auxiliary\_classifier, na qual realiza a classificação auxiliar do modelo; e build\_googlenet, que finalmente realiza a construção do modelo.

\lstinputlisting[
    breaklines=true,
    language={Python},
    caption={Código fonte da construção do modelo GoogleNet},
    label={codigo:googlenet},
    style=mestilo
]{Codigo/googlenet.py}

\paragraph{}Outro aspecto fundamental para o treinamento dos modelos foi o desenvolvimento de calbacks para os modelos. Dessa forma, durante o processo de treinamento, o treinamento era interrompido para: checar se houve melhoria em relação ao modelo anterior, podendo abreviar o treinamento caso não tenha ocorrido nenhuma melhoria; e salvamento automático, possibilitando recuperar o modelo em caso de queda de energia ou interrupção durante o treinamento.

\lstinputlisting[
    breaklines=true,
    language={Python},
    caption={Código fonte do callback de abreviação do treinamento em caso de não haver melhora},
    label={codigo:earlystoppingcallback},
    style=mestilo
]{Codigo/early_stopping_callback.py}

\lstinputlisting[
    breaklines=true,
    language={Python},
    caption={Código fonte do callback de salvamento automático do modelo ResNet50},
    label={codigo:checkpointcallback},
    style=mestilo
]{Codigo/checkpoint_callback.py}

\section{Resultados}

\section{Conclusão}



\bibliographystyle{sbc}
\bibliography{bibliografia}

\end{document}
